\section{Introduction}
Continuous reconstruction of individual finger movements is an important algorithmic task for motor brain--computer interfaces. Unlike discrete gesture recognition, trajectory decoding must recover the timing and shape of flexion throughout both movement and rest. ECoG provides spatially resolved cortical recordings from which movement trajectories can be estimated~\cite{schalk2007}. Public recordings and accompanying analyses have made this problem accessible to reproducible method development~\cite{miller2019}.

High-frequency power and low-frequency voltage describe complementary signal properties. Power alone loses waveform polarity, whereas coarse spectral envelopes can discard frequency detail. A decoder must also capture local transitions and movement-sequence context despite short labeled recordings and heterogeneous electrode coverage. These considerations motivate combining structured spectral features, multiscale convolutions, and recurrent context.

We introduce \emph{RecurFlex}, which concatenates a dense high-frequency Morlet representation with four signed low-frequency bands. A temporal encoder--decoder produces sample-aligned trajectories, and a residual bidirectional gated recurrent unit (BiGRU) augments its compressed bottleneck. Skip connections retain local detail while recurrence operates at reduced temporal resolution. We evaluate this combination through complete test-recording reconstruction and controlled feature experiments.

We evaluate official BCI Competition IV Dataset 4 tests and all nine Miller fingerflex subjects under a chronological two-thirds development and one-third testing split. Figure~\ref{fig:model} summarizes the architecture.

\begin{figure*}[t]
\centering
\resizebox{\textwidth}{!}{%
% Original RecurFlex diagram; feature glyphs are schematic.
% Native width 226.6 mm; the manuscript scales it to \textwidth.
\begingroup%
\definecolor{RFink}{RGB}{43,52,65}%
\definecolor{RFblue}{RGB}{69,115,166}%
\definecolor{RFteal}{RGB}{55,135,126}%
\definecolor{RFpurple}{RGB}{128,99,155}%
\definecolor{RFgold}{RGB}{177,132,62}%
\definecolor{RForange}{RGB}{219,135,48}%
\definecolor{RFline}{RGB}{208,214,223}%
\begin{tikzpicture}[
  x=1mm,y=-1mm,
  font=\sffamily\fontsize{10}{11.5}\selectfont,text=RFink,
  flow/.style={-{Stealth[length=1.6mm,width=1.1mm]},draw=black,line width=0.6pt},
  wire/.style={draw=black,line width=0.6pt},
  skip/.style={-{Stealth[length=1.4mm,width=0.95mm]},draw=black,dashed,line width=0.5pt},
  block/.style={draw=RFline,rounded corners=0.8mm,line width=0.55pt,
    align=center,inner sep=0.7mm},
  small/.style={font=\sffamily\fontsize{9.5}{10.8}\selectfont},
  feature/.style={font=\sffamily\fontsize{8.5}{9.8}\selectfont},
  head/.style={font=\sffamily\bfseries\fontsize{10}{11.5}\selectfont,anchor=west}
]
\path[use as bounding box] (-0.3,-1) rectangle (226.3,71);
\begin{pgfinterruptboundingbox}
\node[head] at (1,2) {(a) Multiband features};
\node[head] at (74,2) {(b) RecurFlex};
\node[head,anchor=center] at (210,2) {(c) Task head};

% (a) Top-to-bottom feature pipeline; LF and HF branches sit side by side.
\node[block,fill=RFblue!4,minimum width=22.5mm,minimum height=10mm,feature]
  (ecog) at (34.5,12) {Preprocessed\\ECoG};
\node[block,fill=RFpurple!6,minimum width=28.5mm,minimum height=13mm]
  (lf) at (16.5,32) {};
\node[block,fill=RFblue!5,minimum width=28.5mm,minimum height=13mm]
  (hf) at (52.5,32) {};
\draw[wire] (ecog.south) -- (34.5,21);
\draw[flow] (34.5,21) -| (lf.north);
\draw[flow] (34.5,21) -| (hf.north);
\foreach \k in {0,...,3}{
  \draw[RFpurple!25,line width=0.15pt] (3.7,{29.03+1.89*\k}) -- (9.37,{29.03+1.89*\k});
  \draw[RFpurple,line width=0.55pt]
    plot[domain=0:1,samples=45,variable=\t]
    ({3.7+5.67*\t},{29.03+1.89*\k-0.585*sin((360+100*\k)*\t+28*\k)});
}
\node[align=center,feature] at (20.7,32) {Signed voltage\\$\delta,\theta,\alpha,\beta$\\0.5--30 Hz};
\foreach \i in {0,...,6}{\foreach \j in {0,...,5}{
  \pgfmathtruncatemacro{\tone}{20+mod(13*\i+23*\j+7*\i*\j,65)}
  \fill[RFblue!\tone!white] ({39.7+0.81*\i},{28.85+0.99*\j}) rectangle ++(0.747,0.927);
}}
\node[align=center,feature] at (56.7,32) {Morlet power\\40 log bands\\40--300 Hz};
\draw[wire] (lf.south) -- (16.5,43) -- (34.5,43);
\draw[wire] (hf.south) -- (52.5,43) -- (34.5,43);
\node[block,draw=RFgold,fill=RFgold!17,minimum width=21.5mm,minimum height=9mm,feature]
  (flatten) at (34.5,52) {Flatten};
\draw[flow] (34.5,43) -- (flatten.north);
\draw[flow,rounded corners=1mm]
  (flatten.south) -- (34.5,68.5) -- (71,68.5) -- (71,15) -- (75.3,15);

% (b) Original module sizes; 8.5-mm vertical level spacing.
% The encoder and decoder mirror about x=135.6 mm.
% The five pre-pooling states E0--E4 supply the skip connections.
\foreach \name/\xx/\yy/\hh/\channels in {
  E0/78.4/15/10/64,E1/89.4/23.5/8/64,E2/100.4/32/6.7/128,
  E3/111.4/40.5/5.3/256,E4/122.4/49/4.4/512}{
  \node[minimum width=6.2mm,minimum height=\hh mm,inner sep=0pt]
    (\name) at (\xx,\yy) {};
  \filldraw[fill=RFblue!23,draw=RFblue,line width=0.55pt]
    ({\xx-3.1},{\yy-\hh/2}) rectangle ({\xx+3.1},{\yy+\hh/2});
  \filldraw[fill=RFblue!42,draw=RFblue,line width=0.5pt]
    ({\xx+3.1},{\yy-\hh/2}) -- ++(1,-1) -- ++(0,\hh) -- ++(-1,1) -- cycle;
  \filldraw[fill=RFblue!12,draw=RFblue,line width=0.5pt]
    ({\xx-3.1},{\yy-\hh/2}) -- ++(1,-1) -- ++(6.2,0) -- ++(-1,1) -- cycle;
  \node[small] at (\xx,\yy) {\channels};
}
\foreach \name/\xx/\yy/\hh/\channels in {
  D1/148.8/49/4.4/512,D2/159.8/40.5/5.3/256,D3/170.8/32/6.7/128,
  D4/181.8/23.5/8/64,D5/192.8/15/10/64}{
  \node[minimum width=6.2mm,minimum height=\hh mm,inner sep=0pt]
    (\name) at (\xx,\yy) {};
  \filldraw[fill=RFteal!22,draw=RFteal,line width=0.55pt]
    ({\xx-3.1},{\yy-\hh/2}) rectangle ({\xx+3.1},{\yy+\hh/2});
  \filldraw[fill=RFteal!40,draw=RFteal,line width=0.5pt]
    ({\xx-3.1},{\yy-\hh/2}) -- ++(-1,-1) -- ++(0,\hh) -- ++(1,1) -- cycle;
  \filldraw[fill=RFteal!10,draw=RFteal,line width=0.5pt]
    ({\xx+3.1},{\yy-\hh/2}) -- ++(-1,-1) -- ++(-6.2,0) -- ++(1,1) -- cycle;
  \node[small] at (\xx,\yy) {\channels};
}
\node[small] at (78.4,6.5) {Stem};
\foreach \a/\b in {E0/E1,E1/E2,E2/E3,E3/E4}{
  \draw[flow] (\a.south) |- (\b.west);
}
\foreach \a/\b in {D1/D2,D2/D3,D3/D4,D4/D5}{
  \draw[flow] (\a.east) -| (\b.south);
}
\foreach \src/\dest/\yy/\res in {
  E0/D5/15/T,E1/D4/23.5/{T/2},E2/D3/32/{T/4},
  E3/D2/40.5/{T/8},E4/D1/49/{T/16}}{
  \draw[skip] ([xshift=1mm]\src.east) -- ([xshift=-1mm]\dest.west);
  \node[fill=white,inner sep=0.5mm,small] at (135.6,\yy) {$\res$};
}
\node[block,draw=RFpurple,fill=RFpurple!12,
  minimum width=18mm,minimum height=10mm,small]
  (context) at (135.6,58.5) {Residual\\BiGRU};
\draw[flow] (E4.south) |- (context.west);
\draw[flow] (context.east) -| (D1.south);
\node[small,text=RFpurple,anchor=south,inner sep=0pt,outer sep=0pt] at (135.6,68.5) {$512\times T/32$};
\node[small,text=RFblue,align=center] at (100.4,58.5) {Encoder\\Conv + Max pooling};
\node[small,text=RFteal,align=center] at (170.8,58.5) {Decoder\\Conv + Upsampling};

% (c) Horizontal pointwise task head: nn.Conv1d(128, 5, 1).
\node[block,draw=RFgold,fill=RFgold!17,minimum width=22mm,minimum height=8mm,small]
  (readout) at (210,15) {Conv1D};
\draw[flow] (D5.east) -- (readout.west);
\node[block,draw=RForange,fill=RForange!4,minimum width=22mm,minimum height=44mm]
  (output) at (210,46.5) {};
\begin{scope}
  \clip[rounded corners=0.8mm] (199,24.5) rectangle (221,68.5);
  \fill[RForange!32] (199,24.5) rectangle (221,30.5);
\end{scope}
\draw[RForange,line width=0.55pt] (199,30.5) -- (221,30.5);
\draw[RForange,rounded corners=0.8mm,line width=0.55pt] (199,24.5) rectangle (221,68.5);
\node[small] at (210,27.5) {Output};
% Schematic flexion: flat intervals, repeated sharp peaks, and return to rest.
% These glyphs are not experimental measurements or model predictions.
\foreach \k/\curvecolor in {0/RFblue,1/RFteal,2/RFpurple,3/RFgold,4/RFink}{
  \draw[RFline,line width=0.15pt] (201,{37.5+6.7*\k}) -- (219,{37.5+6.7*\k});
  \draw[\curvecolor,line width=0.55pt] plot coordinates {
    ({201+18*(0)},{37.5+6.7*\k-0.06*(0.90+0.04*mod(\k,3))})
    ({201+18*(0.06)},{37.5+6.7*\k-0.075*(0.90+0.04*mod(\k,3))})
    ({201+18*(0.12)},{37.5+6.7*\k-0.045*(0.90+0.04*mod(\k,3))})
    ({201+18*(0.16+0.008*\k)},{37.5+6.7*\k-0.06*(0.90+0.04*mod(\k,3))})
    ({201+18*(0.18+0.008*\k)},{37.5+6.7*\k-0.045*(0.90+0.04*mod(\k,3))})
    ({201+18*(0.19+0.008*\k)},{37.5+6.7*\k-4.2*(0.90+0.04*mod(\k,3))})
    ({201+18*(0.2+0.008*\k)},{37.5+6.7*\k-0.18*(0.90+0.04*mod(\k,3))})
    ({201+18*(0.213+0.008*\k)},{37.5+6.7*\k-5.25*(0.90+0.04*mod(\k,3))})
    ({201+18*(0.225+0.008*\k)},{37.5+6.7*\k-0.135*(0.90+0.04*mod(\k,3))})
    ({201+18*(0.238+0.008*\k)},{37.5+6.7*\k-4.95*(0.90+0.04*mod(\k,3))})
    ({201+18*(0.251+0.008*\k)},{37.5+6.7*\k-0.18*(0.90+0.04*mod(\k,3))})
    ({201+18*(0.264+0.008*\k)},{37.5+6.7*\k-5.7*(0.90+0.04*mod(\k,3))})
    ({201+18*(0.277+0.008*\k)},{37.5+6.7*\k-0.135*(0.90+0.04*mod(\k,3))})
    ({201+18*(0.29+0.008*\k)},{37.5+6.7*\k-4.8*(0.90+0.04*mod(\k,3))})
    ({201+18*(0.303+0.008*\k)},{37.5+6.7*\k-0.135*(0.90+0.04*mod(\k,3))})
    ({201+18*(0.316+0.008*\k)},{37.5+6.7*\k-4.05*(0.90+0.04*mod(\k,3))})
    ({201+18*(0.329+0.008*\k)},{37.5+6.7*\k-0.12*(0.90+0.04*mod(\k,3))})
    ({201+18*(0.342+0.008*\k)},{37.5+6.7*\k-4.5*(0.90+0.04*mod(\k,3))})
    ({201+18*(0.355+0.008*\k)},{37.5+6.7*\k-0.105*(0.90+0.04*mod(\k,3))})
    ({201+18*(0.39+0.008*\k)},{37.5+6.7*\k-0.06*(0.90+0.04*mod(\k,3))})
    ({201+18*(0.48+0.008*\k)},{37.5+6.7*\k-0.075*(0.90+0.04*mod(\k,3))})
    ({201+18*(0.57+0.008*\k)},{37.5+6.7*\k-0.075*(0.90+0.04*mod(\k,3))})
    ({201+18*(0.68+0.008*\k)},{37.5+6.7*\k-0.06*(0.90+0.04*mod(\k,3))})
    ({201+18*(0.715+0.008*\k)},{37.5+6.7*\k-0.075*(0.90+0.04*mod(\k,3))})
    ({201+18*(0.728+0.008*\k)},{37.5+6.7*\k-5.55*(0.90+0.04*mod(\k,3))})
    ({201+18*(0.741+0.008*\k)},{37.5+6.7*\k-0.165*(0.90+0.04*mod(\k,3))})
    ({201+18*(0.754+0.008*\k)},{37.5+6.7*\k-4.65*(0.90+0.04*mod(\k,3))})
    ({201+18*(0.767+0.008*\k)},{37.5+6.7*\k-0.12*(0.90+0.04*mod(\k,3))})
    ({201+18*(0.8+0.008*\k)},{37.5+6.7*\k-0.06*(0.90+0.04*mod(\k,3))})
    ({201+18*(0.9)},{37.5+6.7*\k-0.075*(0.90+0.04*mod(\k,3))})
    ({201+18*(1)},{37.5+6.7*\k-0.06*(0.90+0.04*mod(\k,3))})
  };
}
\draw[flow] (readout.south) -- (output.north);
\end{pgfinterruptboundingbox}
\end{tikzpicture}%
\endgroup%
%
}
\caption{RecurFlex architecture. (a) Preprocessed ECoG branches into signed low-frequency voltages and Morlet power; their 100-Hz features are merged and flattened. (b) A mirrored encoder--decoder surrounds a residual BiGRU bottleneck. Dashed paths concatenate encoder features with upsampled decoder features; channel counts are shown before concatenation. (c) A pointwise $1\times1$ Conv1D maps 128 concatenated channels to five trajectories. Feature and output glyphs are schematic.}
\label{fig:model}
\end{figure*}

